% theme: quadratic_equations_speed_applications
\MajorQuestionLesson{速さの利用問題}
\SetRowSpace{5mm}

\TypeQuestionSingle{}{douten1}
\FigProblemBlock{50mm}{fig/test_08_01}{\Text{直角三角形 $ABC$ があり、$\Ang{B}=90^\circ$、$AB=12$cm、$BC=16$cmである。点 $P$ は点 $A$ を出発して辺 $AB$ 上を点 $B$ まで毎秒1cmの速さで動き、点 $Q$ は点 $B$ を出発して辺 $BC$ 上を点 $C$ まで毎秒2cmの速さで動く。2点 $P,Q$ が同時に出発してから $t$ 秒後($0<t\leqq 8$)について、次の問いに答えなさい。}}
{
\QQ{$\Lseg{PB},\ \Lseg{BQ}$ の長さを $t$ を使って表しなさい。}{$PB=12-t,\ BQ=2t$}
\QQ{$\Tri{PBQ}$ の面積を $t$ を使って表しなさい。}{$12t-t^2$ cm$^2$}
\QQ{$\Tri{PBQ}$ の面積が32cm$^2$になるのは、出発してから何秒後か求めなさい。}{$4$秒後、$8$秒後}
}

\TypeQuestionSingle{}{douten2}
\FigProblemBlock{50mm}{fig/lesson_08_douten_02}{\Text{直角三角形 $ABC$ があり、$\Ang{B}=90^\circ$、$AB=12$cm、$BC=16$cmである。点 $P$ は点 $A$ を出発して辺 $AB$ 上を点 $B$ まで毎秒2cmの速さで動き、点 $Q$ は点 $C$ を出発して辺 $BC$ 上を点 $B$ まで毎秒1cmの速さで動く。2点 $P,Q$ が同時に出発してから $t$ 秒後($0<t\leqq 6$)について、次の問いに答えなさい。}}
{
\QQ{$\Lseg{PB},\ \Lseg{BQ}$ の長さを $t$ を使って表しなさい。}{$PB=12-2t,\ BQ=16-t$}
\QQ{$\Tri{PBQ}$ の面積を $t$ を使って表しなさい。}{$t^2-22t+96$ cm$^2$}
\QQ{$\Tri{PBQ}$ の面積が39cm$^2$になるのは、出発してから何秒後か求めなさい。}{$3$秒後}
}

\LessonPageBreak

\TypeQuestionSingle{}{douten3}
\FigProblemBlock{50mm}{fig/lesson_08_douten_03}{\Text{1辺10cmの正方形 $ABCD$ がある。点 $P$ は点 $A$ を出発し、辺 $AB$、辺 $BC$ の順に毎秒2cmの速さで動く。点 $Q$ は点 $D$ を出発して辺 $DC$ 上を点 $C$ まで毎秒1cmの速さで動く。2点 $P,Q$ が同時に出発してから $t$ 秒後($0<t\leqq 10$)について、次の問いに答えなさい。}}
{
\QQ{$0<t\leqq 5$ のとき、$\Tri{DPQ}$ の面積を $t$ を使って表しなさい。}{$5t$ cm$^2$}
\QQ{$5<t\leqq 10$ のとき、$\Tri{DPQ}$ の面積を $t$ を使って表しなさい。}{$10t-t^2$ cm$^2$}
\QQ{$\Tri{DPQ}$ の面積が21cm$^2$になるのは、出発してから何秒後かすべて求めなさい。}{$\dfrac{21}{5}$秒後、$7$秒後}
}

\TypeQuestionSingle{}{douten4}
\FigProblemBlock{50mm}{fig/lesson_08_douten_04}{\Text{長方形 $ABCD$ があり、$AB=16$cm、$BC=8$cmである。点 $P$ は点 $A$ を出発して辺 $AB$ 上を点 $B$ まで毎秒1cmの速さで動き、点 $Q$ は点 $C$ を出発して辺 $CD$ 上を点 $D$ まで毎秒2cmの速さで動く。2点 $P,Q$ が同時に出発してから $t$ 秒後($0<t\leqq 8$)について、次の問いに答えなさい。}}
{
\QQ{$\Lseg{PQ}^2$ を $t$ を使って表しなさい。}{$9t^2-96t+320$}
\QQ{$\Lseg{PQ}=9$cm になるのは、出発してから何秒後かすべて求めなさい。}{$\dfrac{16-\sqrt{17}}{3}$秒後、$\dfrac{16+\sqrt{17}}{3}$秒後}
}
